\section{Introduction}
\label{sec:introduction}

Lists are finite sequences of values that support a wide range of
operations in functional and declarative programming. When combined with
summation, they form the backbone for definitions of sequences,
recurrence, accumulation, and integration in the discrete domain.

Our approach mirrors traditional recursive definitions but is formally
verified using
\href{https://epfl-lara.github.io/stainless/intro.html}{Scala Stainless}
\cite{hamza2019systemfr}, a verification framework for pure Scala programs
that uses formal verification to ensure user-defined functions satisfy
given preconditions, postconditions, and invariants through automated
proofs under all valid inputs.

\begin{quote}
Formal verification is the act of proving or disproving the correctness of
intended algorithms underlying a system with respect to a certain formal
specification or property, using formal methods of mathematics.
--- Wikipedia on Formal Verification \cite{wikipedia2026formalverification}
\end{quote}

This article verifies:

\begin{itemize}
  \item Index and access: tail shift, last element (Section~3).
  \item Slice: recursive, index-range, append consistency (Section~4).
  \item Sum: definition, concatenation, commutativity, positivity
    (Section~5).
  \item Product: definition, concatenation, commutativity, positivity
    (Section~6).
  \item Product divisibility: head, all elements, inserted element
    (Section~7).
  \item Bound and order: lower- and upper-bound propagation through append
    and split (Section~8).
  \item Slice equivalence (Section~9).
  \item Shifted list: period, gap identity, gap translation
    (Section~10).
  \item Rotation: permutation invariants (size, sum, bounds, membership)
    (Section~11).
\end{itemize}

\subsection*{Related work}

Recursive lists, indexed access, and list splitting are long-established
parts of formal libraries. The Rocq/Coq standard list library defines
indexed access, prefix and suffix operations, and proves their
reconstruction law \texttt{firstn n l ++ skipn n l = l}
\cite{rocq2024liststdlib}. Lean's mathematical library also formalizes
list rotation through splitting and concatenation, including reduction of
a rotation index modulo the list length \cite{leanmathlibrotate}.

These prior developments are useful points of contact for the present
work. They show how a mature formal-mathematical setting treats familiar
list structure, while this article develops and verifies the stated
property package for a minimal recursive \texttt{BigInt} implementation
in Scala Stainless. In particular, the article brings structural
operations into the same checked development as product divisibility,
numeric bounds, shifted-list gaps, and rotation invariants. The comparison
is contextual rather than competitive: it locates the Stainless proofs in
the wider body of formal work and makes both the overlap and the scope of
this implementation clear.
